\documentclass[11pt,a4paper]{letter}
\usepackage[utf8]{inputenc}
\usepackage[T1]{fontenc}
\usepackage[french]{babel}
\usepackage[left=2.5cm,right=2.5cm,top=2cm,bottom=2cm]{geometry}
\usepackage{xcolor}

% --- Couleurs Valeo ---
\definecolor{primary}{RGB}{0, 125, 50}

\signature{Loïc Aron MBASSI EWOLO}
\address{Loïc Aron MBASSI EWOLO\\
Cergy, France\\
(+33) 7 59 52 33 98\\
loicaron.mbassiewolo@ensea.fr\\
LinkedIn : loic-mbassi | GitHub : Nameless0l}

\begin{document}

\begin{letter}{Valeo\\
Mobility Tech Center\\
Créteil}

\opening{Madame, Monsieur,}

Élève-ingénieur en 2A à l'ENSEA (double diplôme avec Polytechnique Yaoundé), je candidate au stage R\&D algorithmes dans l'équipe Vital Signs.

Le sujet me parle directement : j'ai travaillé sur du traitement de signaux ECG (proche du PPG) et je comprends les enjeux d'extraction de caractéristiques sur des signaux physiologiques bruités.

Côté développement multi-plateforme, j'ai de l'expérience en \textbf{C/C++} (projet système multiprocessus, STM32) et en \textbf{Python}. Sur GPU, j'ai déployé du YOLOv10 sur Nvidia Jetson pour du temps réel -- je n'ai pas encore codé en CUDA directement mais je comprends le modèle de parallélisation et je suis motivé pour apprendre.

Pour l'embarqué, j'ai développé un modèle TinyML sur STM32 pour de la reconnaissance gestuelle, avec les contraintes mémoire et latence que ça implique. C'est exactement le type de compromis qu'il faut faire pour une "version allégée" comme vous le mentionnez.

Je suis à l'ENSEA à Cergy, donc Créteil c'est accessible. Et la perspective d'un CDI ou d'une thèse ensuite m'intéresse.

Disponible pour un entretien.

\closing{Cordialement,}

\end{letter}
\end{document}
